\section{What the verifiers disagree about}
\label{sec:disagree}

Before any assumption or statistics, the verdicts themselves show where the verifiers part ways
and why. Apart from one comparison with the human audit, nothing in this section uses a human
label.

\paragraph{The verifiers agree on almost every response.}
Of the \NumRows{} responses, \CellAgreeAll{} are credited by all three verifiers and
\CellRejectAll{} by none (Table~\ref{tab:cells}). Only the remaining \RowsAnyDisagree{} get a
different verdict depending on the verifier. Where the three agree, the reader nearly always agrees
with them: the reader verdict is $1$ for \AuditCorrectAgreeAll{} of the \AuditNAgreeAll{} audited
responses in cell $111$ and for \AuditCorrectRejectAll{} of the \AuditNRejectAll{} in cell $000$
(Appendix Figure~\ref{fig:cells}). This is what will later let us share labels between models in
those two cells.

% Generated by scripts/fb_extraction_taxonomy.py -- do not edit by hand.
\begin{table}[t]
\centering
\caption{Why the verifiers disagree: every discordant response, classified from what each verifier's extraction step recorded, with no human labels. \emph{Extraction}: a verifier that gave no credit found no answer. \emph{Judgement}: every verifier that gave no credit extracted an answer. \emph{Runtime}: a verifier timed out. The classes are exclusive.}
\label{tab:taxonomy}
\small
\begin{tabular}{lr}
\toprule
class & responses \\
\midrule
extraction & 891 \\
\quad no boxed answer (boxed verifier finds nothing) & 533 \\
\quad template sentence missing (2024 verifier finds nothing) & 205 \\
\quad template+box & 153 \\
judgement & 268 \\
runtime & 0 \\
\midrule
all discordant & 1{,}159 \\
\bottomrule
\end{tabular}
\end{table}


\paragraph{Most disagreements are failures to find an answer.}
Table~\ref{tab:taxonomy} sorts every discordant response by what each verifier's extraction step
recorded, using fixed rules and no human labels (Appendix~\ref{app:gallery}). In \TaxExtraction{}
of the \RowsAnyDisagree{}, a verifier that gave no credit had found no answer to judge:
\TaxBoxOnly{} responses contain no boxed answer, \TaxTemplateOnly{} do not use the template, and
\TaxTemplateBox{} have neither. In the other \TaxJudgement{}, every verifier
that gave no credit did find an answer, and they disagree about which answer the response gives or
whether it matches the gold answer. No disagreement comes from a verifier timing out. Appendix
Table~\ref{tab:gallery} shows one randomly drawn example of each kind.

\paragraph{Following the format and being right are different things.}
\TemplateCompliant{} of the \NumRows{} responses (\TemplateCompliantPct\%) use the template, in the
sense that the 2024 verifier's pattern finds an answer in them, and across models
the share runs from \TemplateMinPct\% to \TemplateMaxPct\% (Appendix Table~\ref{tab:models}). A
published score mixes the two properties. Of the \CreditedOthersNotEm{} responses that the boxed
verifier or Math-Verify credits and the 2024 verifier does not, \CreditedOthersNotEmNoncompliant{}
skip the template; for those, the 2024 verifier's zero means it could not read the answer.

\paragraph{Neither strict nor lenient extraction is safe.}
The errors run both ways, which is why we do not simply pick a verifier. A strict rule like the
2024 verifier's misses correct answers that skip the template: among the \AuditReaderCorrect{}
audited responses with reader verdict $1$, \AuditReaderCorrectNoncompliant{} skip it, and the 2024
verifier credits \AuditReaderCorrectNoncompliantEm{} of them. A lenient rule like the boxed
verifier's or Math-Verify's also finds answers a reader would not accept: among the
\AuditReaderIncorrectNoncompliant{} audited responses that skip the template and have reader verdict
$0$, the boxed verifier credits \AuditReaderIncorrectNoncompliantSb{} and Math-Verify
\AuditReaderIncorrectNoncompliantMv{}.
The audit is not a random sample of the board, so these are counts, not rates.

\paragraph{A change of wording can flip a verdict.}
A check that needs no reader points the same way. We took \PertN{} responses that all three
verifiers credit and rewrote each in ways that leave the stated answer untouched and that occur in
real responses. Replacing the template's closing words with ``Great job!'', as in the introduction's
example, leaves the other two verifiers crediting every response and the 2024 verifier crediting
\PertSignoffEm{}. Removing the template but keeping the box leaves the 2024 verifier crediting
\PertNoTemplateEm{}. Removing the box but keeping the template leaves the 2024 verifier crediting
every response, the boxed verifier \PertNoBoxSb{}, and Math-Verify \PertNoBoxMv{}. Each verifier
can be moved by a change of form that leaves the answer as it was.

So the verifiers disagree on a small share of responses, mostly because one of them could not read
an answer, and how often that happens depends strongly on the model. Whether it is enough to change
the ranking is the question the next two sections answer.
